\documentclass[11pt,a4paper]{article}

\usepackage[margin=1in]{geometry}
\usepackage{amsmath,amssymb,amsthm,mathtools,bm}
\usepackage{microtype}
\usepackage{booktabs}
\usepackage{enumitem}
\usepackage{graphicx}
\usepackage{xcolor}
\usepackage[hidelinks]{hyperref}
\hypersetup{
  pdftitle={Exact Rate--Distortion Theory for Complex-Projective Born Prediction: Finite-Dimensional Bingham Frontiers, Thermodynamic Coexistence, and Worst-State Capacity},
  pdfauthor={Lluis Eriksson},
  pdfsubject={Quantum information theory; exact finite-dimensional and thermodynamic rate--distortion laws for Born-probability prediction},
  pdfkeywords={Born probability, rate distortion, complex projective space, worst-state channel capacity, Kummer asymptotics, thermodynamic limit}
}
\usepackage[backend=biber,style=numeric,sorting=none,maxbibnames=99]{biblatex}
\addbibresource{references.bib}

\newtheorem{theorem}{Theorem}[section]
\newtheorem{proposition}[theorem]{Proposition}
\newtheorem{lemma}[theorem]{Lemma}
\newtheorem{corollary}[theorem]{Corollary}
\theoremstyle{remark}
\newtheorem{remark}[theorem]{Remark}

\newcommand{\CP}{\mathbb{CP}}
\newcommand{\E}{\mathbb{E}}
\newcommand{\Pp}{\mathbb{P}}
\newcommand{\R}{\mathbb{R}}
\newcommand{\cR}{\mathcal{R}}
\newcommand{\cC}{\mathcal{C}}
\newcommand{\Tr}{\operatorname{tr}}
\newcommand{\KL}{D_{\mathrm{KL}}}
\newcommand{\dd}{\,\mathrm{d}}

\title{\bfseries Exact Rate--Distortion Theory for Complex-Projective Born Prediction:\\
Finite-Dimensional Bingham Frontiers, Thermodynamic Coexistence, and Worst-State Capacity}
\author{Lluis Eriksson}
\date{August 11, 2026}

\begin{document}
\maketitle

\begin{abstract}
We give a unified exact rate--distortion theory for retaining the complete
rank-one Born-probability field of a Haar-random pure state in
$\mathbb{CP}^{d-1}$.  First, a constrained root count, a sharp centered
power-sum theorem, and Newton's expansion prove an all-degree spectral result:
at fixed traceless Frobenius norm the projective Laplace transform is maximized
by a one-positive-spike spectrum.  This reduces the unrestricted Shannon
rate--distortion function to an exact scalar complex-Bingham envelope in every
dimension, with covariant channels and independently flagged coexistence
mixtures attaining every distortion.  For $d\ge3$ directional information
turns on discontinuously; the qubit curve and the all-dimensional
high-fidelity constant are explicit.

We then solve the full fixed-normalized-distortion limit $d\to\infty$.  If
$y_*>2$ solves $y_*-1=2\log y_*$, then the limiting free energy has a unique
coexistence point with
$\alpha_*=y_*^2/[2(y_*-1)]=2.455407482\ldots$ and
$b_*=1-y_*^{-1}=0.715331863\ldots$.  The information cost per dimension is
\[
 r_\infty(\delta)=
 \begin{cases}
 -\log(1-\sqrt{1-\delta}),&0<\delta\le1-b_*^2,\\
 \alpha_*(1-\delta),&1-b_*^2\le\delta\le1,
 \end{cases}
\]
and the finite-dimensional onset multiplier obeys
$\lambda_{c,d}=\alpha_*d-[\alpha_*/(2\log y_*)]\log d+O(1)$.

Finally, a pointwise Hilbert projection converts every measurable scalar
reporter into a physical density-matrix reporter without increasing risk for
any pure input.  Hence the exact worst-state channel capacity equals the Haar
rate--distortion function, arbitrary joint $n$-state memories cost exactly
$n\cR_d(D)$, and transitivity forces zero rate dispersion with an exponential
finite-blocklength strong converse.  These results concern reusable calibrated
probability fields, not click-only simulation, state update, or intrinsic
randomness.
\end{abstract}

\section{Introduction}

Let $x\in\CP^{d-1}$ label a pure state
$\rho_x=|x\rangle\langle x|$.  A classical representation of $x$ is asked
for the Born probability
\begin{equation}
 p_x(y)=\Tr(\rho_x\rho_y)
 \label{eq:born-field}
\end{equation}
of every rank-one query $y$.  The question is not how many outcomes a single
measurement can generate, but how much classical information is required to
retain an entire calibrated probability field to prescribed accuracy.

For a Haar source, integrated squared calibration loss reduces exactly to
squared Frobenius reconstruction of a Haar rank-one projector.  The resulting
complex-projective Shannon rate--distortion problem has two linked layers.
At finite $d$, the reproduction alphabet is the full density-matrix body and
must first be reduced globally to a scalar complex-Bingham envelope.  We prove
that reduction in all dimensions and show that the first directional phase is
discontinuous for every $d\ge3$.  At large $d$, the active Bingham field is of
order $d$, so small-parameter expansions do not apply directly
\cite{BagyanRichards2023}.  The second layer is the global thermodynamic
optimization of that exact finite-dimensional envelope.

Large-parameter
Kummer asymptotics themselves are classical
\cite{Temme1978Kummer,LopezPagola2010Kummer}; our contribution begins with
their global radius optimization, convex conjugate, coexistence law, and
finite-size shift.  A uniform Laplace principle produces the limiting free
energy with coexistence root $y_*-1=2\log y_*$.  General metric-space and
complex-Grassmann results provide bounds or high-resolution dimensions, while
direct complex-projective and Grassmann quantization treats finite codebooks
or metric-ball asymptotics
\cite{MondalDuttaHeath2006,DaiLiuRider2008,PitavalEtAl2016,RieglerKolianderBolcskei2023};
these works do not evaluate the stochastic all-distortion radial envelope.
The exact real-sphere analysis of Dytso and Cardone is the closest
symmetry-reduction template, but its Bessel/von-Mises--Fisher partition and
full-sphere orbit do not reduce to the Beta/Kummer projective orbit
\cite{DytsoCardone2024}.

The centered-projector map embeds $\CP^{d-1}$ as a proper $U(d)$ orbit of real
dimension $2d-2$ in the $(d^2-1)$-dimensional traceless-Hermitian space.  Its
invariant overlap has law $\mathrm{Beta}(1,d-1)$, not the symmetric coordinate
law of a uniform real sphere.  Thus the real-sphere RDF cannot be imported by
setting the sphere dimension to $2d-2$; in particular its limiting convex
geometry has no counterpart of the projective coexistence face.

We also strengthen the operational scope.  We call the resulting
worst-input channel-capacity frontier ``source-universal'' only in this defined
sense; it is not a claim of a single finite-block universal code for an
unknown source distribution.  Compound-source, invariant, and minimax
rate--distortion theories have a substantial prior literature
\cite{Sakrison1969,Harremoes2009CompactGroups,MahmoodWagner2022}.  The new
content here is the pointwise arbitrary-reporter projection and exact
complex-projective evaluation.  The same argument gives an exact equality for
arbitrary joint memories.  Finally, the invariant optimal channel has
constant $D$-tilted information, so general nonasymptotic lossy-source
converses specialize to a zero-dispersion exponential bound
\cite{KostinaVerdu2012}.

Our contributions are:

\begin{enumerate}[leftmargin=2em]
\item an all-degree one-spike projective partition theorem obtained from an
      exactly identified centered power-sum input, a global constrained root
      count, parity closure, and Newton's expansion;
\item the exact unrestricted finite-dimensional RDF, attained by covariant
      complex-Bingham channels and flagged coexistence mixtures;
\item a discontinuous finite-dimensional onset for every $d\ge3$, together
      with the qubit specialization and exact high-fidelity constant;
\item an exact thermodynamic free energy, a closed limiting RDF at every fixed
      normalized distortion, and dimension-independent coexistence constants
      $\alpha_*,b_*,\delta_*$;
\item convergence of first contact with the finite-size correction
      $\lambda_{c,d}=\alpha_*d-c_*\log d+O(1)$;
\item a pointwise projection of arbitrary probability reporters onto physical
      density reports, valid for Haar and exact projective $2$-design queries;
\item an exact worst-state channel-capacity equality and its finite-block
      tensorization; and
\item zero rate dispersion and an exponential strong converse.
\end{enumerate}

We do not claim a unique positive branch at every finite $d$, a matching
third-order coding theorem, or a rank-$r$ Grassmann extension.  The last would
require a new global matrix-Bingham/HCIZ spectral extremum.

This manuscript supersedes the earlier unpublished drafts entitled
\emph{Exact Information--Accuracy Frontiers for Born-Probability Prediction:
Complex-Bingham Rate--Distortion and a Discontinuous Directional-Information
Onset} and \emph{Thermodynamic Limit of Complex-Projective Born-Prediction
Rate--Distortion: Exact Coexistence and a Logarithmic Onset Shift}.  Their
finite-dimensional and thermodynamic results are incorporated here into one
self-contained article; they are not intended as separate overlapping
publications.

\section{Finite-dimensional starting point}

Let $X$ be Haar on $\CP^{d-1}$ and set
\begin{equation}
 S_X=\rho_X-I/d,
 \qquad R_{0,d}^2=\|S_X\|_F^2=1-\frac1d.
 \label{eq:centered-projector}
\end{equation}
For a density-matrix reproduction $\widehat\rho$, use distortion
$\|\rho_X-\widehat\rho\|_F^2$.  Denote the unrestricted classical
rate--distortion function by
\begin{equation}
 \cR_d(D)=\inf I(X;\widehat\rho),
 \qquad \E\|\rho_X-\widehat\rho\|_F^2\le D.
 \label{eq:rdf-definition}
\end{equation}
The infimum ranges over arbitrary standard-Borel stochastic reproductions.

Let
\begin{equation}
 M_d(s)={}_1F_1(1;d;s)
       =(d-1)\int_0^1e^{st}(1-t)^{d-2}\dd t,
 \qquad
 K_d(s)=\log M_d(s)-\frac{s}{d}.
 \label{eq:MdKd}
\end{equation}
The complex-Bingham and complex-Watson families and their normalizers are
established directional-statistics objects
\cite{Kent1994ComplexBingham,MardiaDryden1999}.  Here $K_d$ is the cumulant generating function of
$T_d-1/d$ for $T_d\sim\mathrm{Beta}(1,d-1)$.  Define
\begin{equation}
 G_{d,\lambda}(b)=K_d(2\lambda b)
                  -\lambda R_{0,d}^2b^2,
 \qquad 0\le b\le1.
 \label{eq:finite-free-energy}
\end{equation}
It is convenient to denote the optimized Gibbs log partition by
\begin{equation}
 c_d(\lambda)=-\lambda R_{0,d}^2+\max_{0\le b\le1}G_{d,\lambda}(b).
 \label{eq:cd-definition}
\end{equation}

We use the following exact finite-dimensional theorem as the starting point.
Its nontrivial spectral input is an all-degree Dirichlet partition extremum:
among traceless Hermitian matrices of fixed Frobenius norm, the projective
Laplace transform is maximized by one positive eigenvalue and $d-1$ equal
negative eigenvalues \cite{BrazitikosPandis2025}.

\begin{theorem}[Finite-dimensional projective envelope]
\label{thm:finite-envelope}
For $0\le D\le R_{0,d}^2$,
\begin{equation}
 \cR_d(D)=\sup_{\lambda\ge0}
 \left\{\lambda(R_{0,d}^2-D)
 -\max_{0\le b\le1}G_{d,\lambda}(b)\right\}.
 \label{eq:finite-envelope}
\end{equation}
Every positive active radius is attained by a covariant complex-Bingham
channel whose posterior mean is
$I/d+b(\rho_u-I/d)$.  Independent revealed mixtures of active radii attain
every nonexposed distortion.
\end{theorem}

Appendix~\ref{app:finite} proves the spectral reduction and the complete Gibbs
equality mechanism needed by Theorem~\ref{thm:finite-envelope}, including the
normalization interface to the external centered power-sum theorem.

\section{Finite-dimensional phase structure and endpoint regimes}
\label{sec:finite-phase}

The exact envelope already implies a qualitative transition that is absent in
the qubit case.  The active radii are precisely the maximizers of
$G_{d,\lambda}$.

\begin{theorem}[Discontinuous finite-dimensional onset]
\label{thm:finite-first-order}
For every $d\ge3$, there exists
\begin{equation}
 0<\lambda_{c,d}<\lambda_{0,d}:=\frac{d(d+1)}2
 \label{eq:finite-lambdac}
\end{equation}
such that $b=0$ and at least one $b_{c,d}>0$ are global maximizers of
$G_{d,\lambda_{c,d}}$.  The positive maximizer cannot approach zero at first
contact.  Consequently the exact frontier has a nontrivial linear face
adjacent to $(D,R)=(R_{0,d}^2,0)$, and optimal directional information turns
on discontinuously.
\end{theorem}

\begin{proof}
Put $W=T_d-1/d$.  It is centered and lies in an interval of length one.
Hoeffding's lemma gives $K_d(s)\le s^2/8$ for every real $s$.  Hence,
uniformly in $b\in[0,1]$,
\begin{equation}
 G_{d,\lambda}(b)\le
 \lambda b^2\left(\frac{\lambda}{2}-R_{0,d}^2\right)<0
 \qquad\left(b>0,\ 0<\lambda<2R_{0,d}^2\right).
 \label{eq:small-lambda}
\end{equation}
Thus $b=0$ is initially the unique global maximizer and
$c_d(\lambda)=-\lambda R_{0,d}^2$ on a nonzero interval.

The variance and third centered moment of
$T_d\sim\operatorname{Beta}(1,d-1)$ are
\begin{equation}
 v_d=\frac{d-1}{d^2(d+1)},\qquad
 \mu_{3,d}=\frac{2(d-1)(d-2)}{d^3(d+1)(d+2)}.
 \label{eq:finite-cumulants}
\end{equation}
Expansion at $b=0$ gives
\begin{equation}
 G_{d,\lambda}(b)=
 \bigl(-\lambda R_{0,d}^2+2\lambda^2v_d\bigr)b^2
 +\frac43\lambda^3\mu_{3,d}b^3+O(b^4).
 \label{eq:finite-landau}
\end{equation}
The quadratic coefficient vanishes at $\lambda_{0,d}=d(d+1)/2$.  For
$d\ge3$, $\mu_{3,d}>0$, so a sufficiently small positive $b$ has positive
gain at $\lambda_{0,d}$, and the same remains true just below it.

Let $\lambda_{c,d}$ be the infimum of multipliers whose global argmax contains
a positive radius.  Estimate~\eqref{eq:small-lambda} makes it positive, and
the preceding expansion makes it strictly smaller than $\lambda_{0,d}$.
Choose positive active radii $b_n$ at
$\lambda_n\downarrow\lambda_{c,d}$.  Since
$K_d(s)/s^2\to v_d/2$ as $s\to0$ and
$\lambda_{c,d}<\lambda_{0,d}$, there exist $\varepsilon,\eta>0$ such that
$G_{d,\lambda}(b)\le-\eta b^2$ for
$0<b\le\varepsilon$ and $\lambda$ near $\lambda_{c,d}$.  Thus every $b_n$ is
at least $\varepsilon$.  Compactness and continuity yield a limiting
$b_{c,d}\ge\varepsilon$ coexisting globally with $0$.  Finally,
Danskin's relation~\eqref{eq:appendix-danskin} and an independently revealed
mixture of the two active channels give the linear face.
\end{proof}

\begin{remark}
The theorem proves first contact and its nonzero jump.  It does not assert
uniqueness of the positive maximizer for every larger multiplier or exclude a
later branch exchange.  The exact envelope and flagged mixtures already cover
such possibilities.
\end{remark}

\begin{figure}[t]
 \centering
 \includegraphics[width=\textwidth]{../repro/scalar_envelope.pdf}
 \caption{Numerical evaluation of the exact finite-dimensional scalar
 envelope for $d=2,3,4,5$.  Left: $\cR_d(D)$ versus normalized distortion.
 Right: one selected globally active radius versus scaled multiplier.  The
 qubit radius emerges continuously, while the higher-dimensional curves show
 the finite jump proved in Theorem~\ref{thm:finite-first-order}.  The bounded
 scalar computation is a reproducibility diagnostic, not a proof input.}
 \label{fig:finite-envelope}
\end{figure}

For $d=2$, the Bloch representation maps centered projectors to the real unit
two-sphere up to a factor $1/2$.  The exact result is therefore the $n=3$
real-sphere theorem of Dytso and Cardone \cite{DytsoCardone2024}.  Put
$r=\sqrt{1-2D}$ and choose $\eta\ge0$ from
\begin{equation}
 r=\coth\eta-\frac1\eta.
 \label{eq:qubit-langevin}
\end{equation}
For $0<D<1/2$, such an $\eta$ exists uniquely and
\begin{equation}
 \cR_2(D)=\eta r-\log\frac{\sinh\eta}{\eta}.
 \label{eq:qubit-rdf}
\end{equation}
Moreover $\cR_2(0)=+\infty$ by the limit $\eta\to\infty$, while
$\cR_2(D)=0$ for $D\ge1/2$.  This is a consistency check and operational
specialization, not a new real-sphere theorem.

\begin{proposition}[Exact fixed-dimensional high-fidelity asymptotic]
\label{prop:high-rate}
For every fixed $d\ge2$,
\begin{equation}
 \cR_d(D)=
 (d-1)\log\frac{2(d-1)}{D}
 -(d-1)-\log\Gamma(d)+o(1),
 \qquad D\downarrow0.
 \label{eq:high-rate}
\end{equation}
\end{proposition}

\begin{proof}
Put $m=d-1$.  For integer $d$,
\begin{equation}
 M_d(s)=\Gamma(d)s^{-m}
 \left(e^s-\sum_{j=0}^{m-1}\frac{s^j}{j!}\right).
 \label{eq:M-integer}
\end{equation}
For $s>0$ let
$\ell_d(s)=\log(1-e^{-s}\sum_{j=0}^{m-1}s^j/j!)\le0$.
For a traceless Hermitian report $A$, define
\begin{equation}
 Q_\lambda(A)=\int
 \exp\{-\lambda\|S_x-A\|_F^2\}\,d\mu_d(x).
 \label{eq:high-rate-partition}
\end{equation}
For the axial reproduction $A_{b,u}=b(\rho_u-I/d)$,
\begin{equation}
 \log Q_\lambda(A_{b,u})=
 -m\log(2\lambda)+\log\Gamma(d)
 -\lambda R_{0,d}^2(1-b)^2-m\log b+\ell_d(2\lambda b).
 \label{eq:high-rate-identity}
\end{equation}
Taking $b=1$ gives a lower bound with vanishing error.  For the upper bound,
$K_d(s)\le R_{0,d}^2s$ for $s\ge0$, so radii $b\le1/2$ are eventually
excluded because their radial gain is at most
$-\lambda R_{0,d}^2/4$.  On $b\ge1/2$, use $-\log b\le2(1-b)$ and
$\ell_d\le0$ to bound the remaining correction by
\begin{equation}
 -\lambda R_{0,d}^2(1-b)^2+2m(1-b)
 \le\frac{m^2}{\lambda R_{0,d}^2}=o(1).
\end{equation}
Consequently
\begin{equation}
 c_d(\lambda)=-m\log(2\lambda)+\log\Gamma(d)+o(1).
 \label{eq:cd-high-rate}
\end{equation}
Secant inequalities for the convex function $c_d$ imply
$-\lambda g_\lambda\to m$ for every
$g_\lambda\in\partial c_d(\lambda)$.  If $\lambda_D$ maximizes the exact
conjugate, then $-D\in\partial c_d(\lambda_D)$ and
$\lambda_DD\to m$.  Therefore
\begin{align}
 \cR_d(D)
 &=-\lambda_DD-c_d(\lambda_D)\\
 &=-m+m\log(2\lambda_D)-\log\Gamma(d)+o(1)\\
 &=m\log\frac{2m}{D}-m-\log\Gamma(d)+o(1),
\end{align}
which proves \eqref{eq:high-rate}.
\end{proof}

The limits $d\to\infty$ at fixed normalized distortion and
$D\downarrow0$ at fixed $d$ are distinct.  The following sections resolve the
first regime; a uniform double-scaling theory remains open.

\section{Uniform projective Laplace principle}

The transition occurs at $\lambda$ of order $d$.  We first isolate the
corresponding asymptotic partition.

\begin{lemma}[Uniform beta Laplace principle]
\label{lem:laplace}
For $y\ge0$, put
\begin{equation}
 F_d(y)=\frac1d\log M_d(dy).
\end{equation}
Then $F_d\to\Phi$ locally uniformly on $[0,\infty)$, where
\begin{equation}
 \Phi(y)=
 \begin{cases}
 0,&0\le y\le1,\\
 y-1-\log y,&y>1.
 \end{cases}
 \label{eq:Phi}
\end{equation}
\end{lemma}

\begin{proof}
Rewrite the beta integral as
\begin{equation}
  M_d(dy)=(d-1)\int_0^1
  \exp\{dyt+(d-2)\log(1-t)\}\dd t.
\end{equation}
For fixed $\varepsilon>0$, the usual compact-interval Laplace bounds apply on
$[0,1-\varepsilon]$.  On the remaining interval the original factor
$(1-t)^{d-2}$ supplies exponential control; letting $\varepsilon$ decrease
after $d\to\infty$ gives the pointwise limit
\begin{equation}
 \sup_{0\le t<1}\{yt+\log(1-t)\}.
\end{equation}
The maximizer is $t=0$ for $y\le1$ and $t=1-y^{-1}$ for $y>1$, which yields
\eqref{eq:Phi}.  Moreover,
\begin{equation}
 F_d'(y)=\E_{d,y}T_d\in[0,1],
\end{equation}
where $\E_{d,y}$ is expectation under exponential tilt $dy$.  Hence the
family is equi-Lipschitz.  A finite-net argument upgrades pointwise
convergence to local-uniform convergence because $\Phi$ is continuous.
\end{proof}

For $\alpha\ge0$, define the normalized optimized gain
\begin{equation}
 g_d(\alpha)=\frac1d\max_{0\le b\le1}G_{d,\alpha d}(b).
 \label{eq:gd}
\end{equation}

\begin{corollary}[Thermodynamic free-energy convergence]
\label{cor:gd}
The functions $g_d$ converge locally uniformly to
\begin{equation}
 g(\alpha)=\max_{0\le b\le1}
 \{\Phi(2\alpha b)-\alpha b^2\}.
 \label{eq:g-limit}
\end{equation}
\end{corollary}

\begin{proof}
Lemma~\ref{lem:laplace}, $R_{0,d}^2\to1$, and
$d^{-1}K_d(2\alpha db)=F_d(2\alpha b)-2\alpha b/d$ give uniform convergence
of the objective on compact $\alpha$ intervals and $b\in[0,1]$.  Taking a
maximum preserves it.
\end{proof}

\section{Global coexistence and the dimension-independent constants}

The limiting scalar problem is globally solvable.

\begin{theorem}[Exact limiting phase diagram]
\label{thm:phase}
Let $y_*>2$ be the unique nontrivial root of
\begin{equation}
 y_*-1=2\log y_*.
 \label{eq:ystar}
\end{equation}
Define
\begin{equation}
 \alpha_*=\frac{y_*^2}{2(y_*-1)},
 \qquad b_*=1-\frac1{y_*}.
 \label{eq:constants}
\end{equation}
Then $g(\alpha)=0$ for $0\le\alpha\le\alpha_*$.  At $\alpha_*$ the global
maximizers in \eqref{eq:g-limit} are exactly $0$ and $b_*$.  For
$\alpha>\alpha_*$ the positive global maximizer is unique and equals
\begin{equation}
 b_+(\alpha)=1-\frac1{y_+(\alpha)},
 \qquad
 y_+(\alpha)=\alpha+\sqrt{\alpha(\alpha-2)}.
 \label{eq:positive-branch}
\end{equation}
\end{theorem}

\begin{proof}
Set $y=2\alpha b$.  The optimization becomes
\begin{equation}
 \max_{0\le y\le2\alpha}
 \left\{\Phi(y)-\frac{y^2}{4\alpha}\right\}.
 \label{eq:y-optimization}
\end{equation}
For $y\le1$ the objective is nonpositive and vanishes only at $y=0$.  For
$y>1$, an interior stationary point obeys
\begin{equation}
 1-\frac1y-\frac{y}{2\alpha}=0,
 \qquad
 \alpha=\frac{y^2}{2(y-1)}.
 \label{eq:stationary}
\end{equation}
The two positive stationary points appear at $\alpha=2$; the larger root is
$y_+(\alpha)$ and is the local maximum.  At any stationary point the gain is
\begin{equation}
 \Phi(y)-\frac{y^2}{4\alpha}
 =\frac{y-1}{2}-\log y.
 \label{eq:stationary-gain}
\end{equation}
The function $y-1-2\log y$ decreases on $(1,2)$, increases on
$(2,\infty)$, and has exactly the roots $1$ and $y_*$.  Thus the larger
stationary point first ties the zero-radius solution at $y_*$.  More
explicitly, for $y>1$,
\begin{equation}
 \frac{\dd}{\dd y}\left(\Phi(y)-\frac{y^2}{4\alpha}\right)
 =-\frac{(y-y_-)(y-y_+)}{2\alpha y}.
\end{equation}
For $\alpha<2$ the derivative is strictly negative.  For $\alpha\ge2$,
$y_-$ is a local minimum and $y_+$ is the only local maximum; the objective
decreases after $y_+$, so the endpoint $y=2\alpha$ cannot improve it.
Equation~\eqref{eq:stationary-gain} then decides exactly when that maximum
beats zero and proves the complete statement.
\end{proof}

Numerically,
\begin{equation}
\begin{split}
 y_*&=3.512862417252\ldots,\\
 \alpha_*&=2.455407482284\ldots,\\
 b_*&=0.715331862959\ldots,\\
 \delta_*:=1-b_*^2&=0.488300325835\ldots .
\end{split}
\label{eq:numerical-constants}
\end{equation}

\section{Closed thermodynamic information--accuracy frontier}

We normalize distortion by the zero-information value $R_{0,d}^2$.

\begin{theorem}[Thermodynamic Born-prediction frontier]
\label{thm:thermodynamic-rdf}
For each fixed $0<\delta\le1$,
\begin{equation}
 \lim_{d\to\infty}\frac1d
 \cR_d(\delta R_{0,d}^2)=r_\infty(\delta),
 \label{eq:rdf-limit}
\end{equation}
where
\begin{equation}
 r_\infty(\delta)=
 \begin{cases}
 -\log(1-\sqrt{1-\delta}),
   &0<\delta\le\delta_*,\\[1mm]
 \alpha_*(1-\delta),
   &\delta_*\le\delta\le1.
 \end{cases}
 \label{eq:rinfty}
\end{equation}
The branches agree at $\delta_*$, where
$r_\infty(\delta_*)=\log y_*$.
\end{theorem}

\begin{proof}
Theorem~\ref{thm:finite-envelope}, with $\lambda=\alpha d$, gives
\begin{equation}
 \frac1d\cR_d(\delta R_{0,d}^2)
 =\sup_{\alpha\ge0}
 \{\alpha R_{0,d}^2(1-\delta)-g_d(\alpha)\}.
 \label{eq:scaled-dual}
\end{equation}
We first justify passage of the supremum to the limit.  Fix $\delta>0$ and
put $\eta=\delta/8$.  The beta tail is exact:
$\Pp\{T_d\ge1-\eta\}=\eta^{d-1}$.  Evaluating the partition on this event
and choosing $b=1$ gives
\begin{equation}
 g_d(\alpha)\ge
 \alpha\left(1-2\eta-\frac1d\right)
 +\left(1-\frac1d\right)\log\eta.
\end{equation}
Consequently, for every $d\ge2$ and $\alpha\ge0$,
\begin{equation}
 \alpha R_{0,d}^2(1-\delta)-g_d(\alpha)
 \le -\frac{\delta}{4}\alpha
 -\left(1-\frac1d\right)\log\eta.
 \label{eq:equicoercive}
\end{equation}
Thus all maximizing multipliers lie in a common compact interval.  Local
uniform convergence from Corollary~\ref{cor:gd} now yields
\begin{equation}
 r_\infty(\delta)=\sup_{\alpha\ge0}
 \{\alpha(1-\delta)-g(\alpha)\}.
 \label{eq:limit-conjugate}
\end{equation}

The right derivative of $g$ jumps at $\alpha_*$ from $0$ to $b_*^2$.
Therefore gains $1-\delta\le b_*^2$ lie on the coexistence face and give
$\alpha_*(1-\delta)$.  If $1-\delta>b_*^2$, envelope differentiation on the
unique positive branch gives
\begin{equation}
 1-\delta=g'(\alpha)=b_+(\alpha)^2.
\end{equation}
Put $b=\sqrt{1-\delta}$ and $y=(1-b)^{-1}$.  Substitution into the conjugate
is explicit: at a stationary point
\begin{equation}
 \alpha b^2-g(\alpha)=2\alpha b^2-\Phi(y)
 =y-1-\Phi(y)=\log y.
\end{equation}
Thus it reduces to $\log y=-\log(1-\sqrt{1-\delta})$.  At coexistence,
\eqref{eq:ystar} makes this value $\log y_*=\alpha_*b_*^2$.
\end{proof}

The fixed-$\delta$ order of limits is essential.  The theorem does not assert
uniformity as $\delta\downarrow0$; the fixed-dimensional high-resolution
limit is a different asymptotic regime.

\begin{figure}[t]
 \centering
 \includegraphics[width=\textwidth]{../repro/thermodynamic_frontier.pdf}
 \caption{Left: the exact limiting information--accuracy curve and its
 coexistence point $\delta_*$.  Right: finite-dimensional first-contact
 multipliers from the exact scalar problem, the leading logarithmic
 approximation, and the limit $\alpha_*$.  The plotted finite-dimensional
 values are reproducibility diagnostics, not proof inputs.}
 \label{fig:thermodynamic}
\end{figure}

\section{Thermodynamic first contact and its logarithmic finite-size correction}

Retain the first multiplier $\lambda_{c,d}$ at which a positive radius is a
global maximizer of \eqref{eq:finite-free-energy}, and choose any positive
active radius $b_{c,d}$ there.

\begin{lemma}[Beta-mgf exclusion bound]
\label{lem:mgf-bound}
For $0\le s<d$,
\begin{equation}
 K_d(s)\le-\log(1-s/d)-s/d.
 \label{eq:mgf-bound}
\end{equation}
Consequently a minimizing first-contact tilt cannot satisfy $s/d\to0$.
\end{lemma}

\begin{proof}
The hypergeometric series and $(d)_k\ge d^k$ give
\begin{equation}
 M_d(s)=\sum_{k\ge0}\frac{s^k}{(d)_k}
 \le\sum_{k\ge0}(s/d)^k=(1-s/d)^{-1}.
\end{equation}
After centering, this is \eqref{eq:mgf-bound}.  If $u=s/d\to0$, then the
right side is $u^2/2+O(u^3)$, so the ratio term is at least
$(1+o(1))R_{0,d}^2d^2/2$.  On the other hand, evaluating
\eqref{eq:contact-infimum} at $s=dy_*$ gives $\lambda_{c,d}=O(d)$.
Therefore a minimizing tilt cannot have $s/d\to0$.
\end{proof}

\begin{lemma}[Uniform interior Kummer--Laplace expansion]
\label{lem:uniform-saddle}
Let $J\Subset(1,\infty)$ be compact.  Uniformly for $y\in J$,
\begin{equation}
 M_d(dy)=e^{d\Phi(y)}y\sqrt{2\pi d}\,[1+O_J(d^{-1})],
 \label{eq:uniform-saddle-prefactor}
\end{equation}
and hence
\begin{equation}
 \log M_d(dy)=d\Phi(y)+\frac12\log d
 +\log(y\sqrt{2\pi})+O_J(d^{-1}).
 \label{eq:uniform-laplace-expansion}
\end{equation}
\end{lemma}

\begin{proof}
Write
\begin{equation}
 M_d(dy)=(d-1)\int_0^1 e^{d f_y(t)}g(t)\dd t,
 \qquad f_y(t)=yt+\log(1-t),\quad g(t)=(1-t)^{-2}.
\end{equation}
The unique maximizer is $t_y=1-y^{-1}$, and as $y$ ranges over $J$ the
points $t_y$ remain in a fixed compact subinterval of $(0,1)$.  Choose
$\rho>0$ so that $[t_y-\rho,t_y+\rho]\subset(0,1)$ for every $y\in J$.
Compactness and strict concavity give a uniform $c_J>0$ such that
\begin{equation}
 f_y(t)\le\Phi(y)-c_J
 \quad\text{whenever }|t-t_y|\ge\rho, y\in J.
 \label{eq:uniform-laplace-tail}
\end{equation}
Thus the exterior integral is exponentially smaller than the main term.

On the interior neighborhood,
\begin{equation}
 f_y''(t_y)=-y^2,\qquad f_y'''(t_y)=-2y^3,\qquad
 f_y^{(4)}(t_y)=-6y^4,\qquad g(t_y)=y^2.
\end{equation}
All derivatives of $f_y$ through order five and of $g$ through order three
are uniformly bounded there.  With $z=\sqrt d\,(t-t_y)$, Taylor's theorem
through fourth order gives a Gaussian integrand
$e^{-y^2z^2/2}$ times a factor whose even contribution is
$1+O_J(d^{-1})(1+|z|^6)$; the order-$d^{-1/2}$ term is odd and integrates to
zero on the symmetric neighborhood.  On
$|z|\le d^{1/10}$ the remainder is uniformly dominated by an integrable
Gaussian times $O_J(d^{-1})(1+|z|^{10})$; on the remainder of the interior
neighborhood strict concavity gives an exponentially small bound.  Therefore
\begin{align}
 M_d(dy)
 &=(d-1)e^{d\Phi(y)}g(t_y)
   \sqrt{\frac{2\pi}{d|f_y''(t_y)|}}\,[1+O_J(d^{-1})]\\
 &=(d-1)e^{d\Phi(y)}y^2
   \sqrt{\frac{2\pi}{dy^2}}\,[1+O_J(d^{-1})]\\
 &=e^{d\Phi(y)}y\sqrt{2\pi d}\,[1+O_J(d^{-1})].
\end{align}
Taking logarithms proves \eqref{eq:uniform-laplace-expansion}.
\end{proof}

\begin{theorem}[First-contact asymptotics]
\label{thm:first-contact}
As $d\to\infty$,
\begin{equation}
 \frac{\lambda_{c,d}}d\to\alpha_*,
 \qquad b_{c,d}\to b_*.
 \label{eq:first-contact-limit}
\end{equation}
The coexistence distortion and rate satisfy
\begin{equation}
 \frac{D_{c,d}}{R_{0,d}^2}\to\delta_*,
 \qquad \frac{R_{c,d}}d\to\log y_*.
 \label{eq:contact-rate-distortion}
\end{equation}
More sharply,
\begin{equation}
 \boxed{\lambda_{c,d}=\alpha_*d-c_*\log d+O(1)},
 \qquad
 c_*=\frac{\alpha_*}{2\log y_*}
    =\frac{y_*^2}{2(y_*-1)^2}.
 \label{eq:log-correction}
\end{equation}
\end{theorem}

\begin{proof}
Write $s=2\lambda b$.  Nonnegative radial gain and $b\le1$ are equivalent to
\begin{equation}
 \lambda\ge \frac{R_{0,d}^2s^2}{4K_d(s)},
 \qquad \lambda\ge\frac{s}{2}.
\end{equation}
Hence
\begin{equation}
 \lambda_{c,d}=\inf_{s>0}
 \max\left\{\frac{s}{2},
 \frac{R_{0,d}^2s^2}{4K_d(s)}\right\}.
 \label{eq:contact-infimum}
\end{equation}
Evaluating this expression at $s=dy_*$ and using local-uniform Laplace
convergence gives $\limsup_d\lambda_{c,d}/d\le\alpha_*$.  This comparison
will be used below to localize every minimizing sequence.
For $s=dy$ with $y>1$, Lemma~\ref{lem:laplace} makes the second term divided
by $d$ converge locally uniformly to
\begin{equation}
 A(y)=\frac{y^2}{4\Phi(y)}.
\end{equation}
Its unique global minimizer is $y_*$: the equation $A'(y)=0$ is exactly
\eqref{eq:ystar}, $A(y_*)=\alpha_*$, and $y_*/2<\alpha_*$.  Lemma
\ref{lem:mgf-bound} excludes $s/d\to0$.  If instead
$s/d\to y\in(0,1]$, local-uniform Laplace convergence gives
$K_d(s)/d\to\Phi(y)=0$, so the ratio term in
\eqref{eq:contact-infimum}, after division by $d$, diverges.  The term $s/2$
excludes $s/d\to\infty$.  Thus every minimizing sequence lies in a compact
subset of $(1,\infty)$, where uniform convergence and uniqueness imply
$s_{c,d}/d\to y_*$ and $\lambda_{c,d}/d\to\alpha_*$.  For all large $d$ the
second term in \eqref{eq:contact-infimum} is active, so
\begin{equation}
 b_{c,d}=\frac{s_{c,d}}{2\lambda_{c,d}}
 =\frac{2K_d(s_{c,d})}{R_{0,d}^2s_{c,d}}
 \longrightarrow\frac{2\Phi(y_*)}{y_*}=b_*.
\end{equation}
The distortion and rate limits follow from
$D=R_{0,d}^2(1-b^2)$ and $R=\lambda R_{0,d}^2b^2$ at coexistence.

It remains to refine the multiplier.  Apply
Lemma~\ref{lem:uniform-saddle} on a compact neighborhood of $y_*$.
Using $K_d(dy)=\log M_d(dy)-y$ and $R_{0,d}^2=1-d^{-1}$ gives
\begin{equation}
 \frac{R_{0,d}^2d^2y^2}{4K_d(dy)}
 =dA(y)-\frac{A(y)}{2\Phi(y)}\log d+O(1)
 \label{eq:contact-expansion}
\end{equation}
uniformly there.  Put $q(y)=A(y)/(2\Phi(y))$.  On a fixed neighborhood $U$
of $y_*$ there are constants $c,L,C>0$ such that
$A(y)\ge A(y_*)+c|y-y_*|^2$, $|q(y)-q(y_*)|\le L|y-y_*|$, and the remainder
in \eqref{eq:contact-expansion} has absolute value at most $C$.  Thus, with
$z=y-y_*$, the right side is bounded below by
\begin{equation}
 dA(y_*)-q(y_*)\log d+dc z^2-L(\log d)|z|-C.
\end{equation}
The variable terms are at least
$-L^2(\log d)^2/(4cd)=o(1)$.  Evaluation at $y=y_*$ gives the matching upper
bound up to $O(1)$.  Since the global minimizer has already been shown to
enter $U$, we obtain
$\lambda_{c,d}=dA(y_*)-q(y_*)\log d+O(1)$.  Finally
$A(y_*)=\alpha_*$ and $\Phi(y_*)=\log y_*$ prove
\eqref{eq:log-correction}.
\end{proof}

\section{Pointwise physical projection and source universality}

So far the source is Haar.  We now allow every pure input state and charge a
representation channel by its Shannon capacity over all input priors.

Let $\nu$ be Haar measure on rank-one queries, or a weighted exact complex
projective $2$-design.  For a scalar report $q\in L^2(\nu)$, define its risk
against $x$ by
\begin{equation}
 \ell_x(q)=\int |q(y)-\Tr(\rho_x\rho_y)|^2\dd\nu(y).
 \label{eq:pointwise-risk}
\end{equation}

\begin{theorem}[Pointwise projection to a physical Born reporter]
\label{thm:pointwise-projection}
For every $q\in L^2(\nu)$ there exists a density matrix $\sigma(q)$ such that
\begin{equation}
 \ell_x\bigl(y\mapsto\Tr(\sigma(q)\rho_y)\bigr)
 \le\ell_x(q)
 \label{eq:pointwise-dominance}
\end{equation}
simultaneously for every pure $x$.  The map $q\mapsto\sigma(q)$ may be chosen
Borel measurable.  Moreover,
\begin{equation}
 \int|\Tr[(\rho_x-\sigma)\rho_y]|^2\dd\nu(y)
 =\frac{\|\rho_x-\sigma\|_F^2}{d(d+1)}.
 \label{eq:design-isometry}
\end{equation}
\end{theorem}

\begin{proof}
Let
\begin{equation}
 \mathcal B=\{y\mapsto\Tr(A\rho_y):A=A^*\}
 \subset L^2(\nu).
\end{equation}
The $2$-design identity makes this a finite-dimensional closed subspace with
matrix norm
\begin{equation}
 \|A\|_*^2=
 \frac{\Tr(A^2)+(\Tr A)^2}{d(d+1)}.
 \label{eq:star-norm}
\end{equation}
Let $f_A$ be the orthogonal projection of $q$ onto $\mathcal B$.  Because
$f_{\rho_x}\in\mathcal B$,
\begin{equation}
 \|f_{\rho_x}-q\|_2^2
 =\|f_{\rho_x}-f_A\|_2^2+\|q-f_A\|_2^2.
\end{equation}
Next let $\sigma$ be the metric projection of $A$ in the norm
\eqref{eq:star-norm} onto the compact convex density-matrix body.  Since
$\rho_x$ belongs to that body, the Hilbert projection inequality gives
\begin{equation}
 \|f_{\rho_x}-f_\sigma\|_2
 \le\|f_{\rho_x}-f_A\|_2.
\end{equation}
This proves \eqref{eq:pointwise-dominance}.  Orthogonal projection onto a
finite-dimensional subspace is continuous in $L^2$, and metric projection
onto a closed convex set in finite dimension is continuous, establishing the
measurable choice.  Finally, both matrices in \eqref{eq:design-isometry} have
trace one, so the trace term in \eqref{eq:star-norm} vanishes.
\end{proof}

Consider a channel $W$ from pure states to an arbitrary standard-Borel
classical memory $Z$, followed by jointly measurable probability reports
$q_Z(y)\in[0,1]$; bounded joint measurability makes
$z\mapsto q_z$ a Borel $L^2(\nu)$-valued map.  Its worst-state
Born risk is $\sup_x\E[\ell_x(q_Z)|X=x]$.  Its information cost is its channel
capacity
\begin{equation}
 C(W)=\sup_P I_P(X;Z),
 \label{eq:channel-capacity}
\end{equation}
where the supremum ranges over all pure-state priors.
For $0\le D\le R_{0,d}^2$, let
\begin{equation}
 \mathfrak W_d(D)=\left\{(W,q):
 \sup_x\E[\ell_x(q_Z)|X=x]\le\frac{D}{d(d+1)}\right\}.
 \label{eq:uniform-risk-class}
\end{equation}
Define the source-universal frontier by
\begin{equation}
 \cC_d^{\rm univ}(D)=
 \inf_{(W,q)\in\mathfrak W_d(D)}\sup_P I_P(X;Z).
 \label{eq:universal-definition}
\end{equation}

\begin{theorem}[Exact source-universal minimax capacity]
\label{thm:minimax}
For every $0\le D\le R_{0,d}^2$, the minimum channel capacity in
\eqref{eq:universal-definition} is exactly
\begin{equation}
 \boxed{\cC_d^{\rm univ}(D)=\cR_d(D)}.
 \label{eq:minimax-equality}
\end{equation}
The Haar prior is least favorable for an optimal covariant channel.
\end{theorem}

\begin{proof}
Theorem~\ref{thm:pointwise-projection} lets us replace every report by a
density report without increasing any statewise risk or channel capacity.
Every uniformly feasible channel is therefore feasible under Haar input, so
\begin{equation}
 C(W)\ge I_{\rm Haar}(X;Z)\ge\cR_d(D).
\end{equation}

Conversely, Theorem~\ref{thm:finite-envelope} supplies a covariant
complex-Bingham channel $W_D$, or an independently flagged mixture of active
channels on a coexistence face, with constant statewise distortion $D$.
Let $Q_D$ be its output law under Haar input.  For one active orbit,
covariance makes the divergence below independent of $x$, and Haar averaging
identifies its value with the attained mutual information.  On a coexistence
face, write the revealed-flag channel and output law as
$W_D=\sum_j a_j\delta_j\otimes W_j$ and
$Q_D=\sum_j a_j\delta_j\otimes Q_j$.  The flag is independent of $x$, so the
same conclusion follows from
$\KL(W_D(\cdot|x)\|Q_D)=\sum_j a_j\KL(W_j(\cdot|x)\|Q_j)$.  Thus
\begin{equation}
 \KL(W_D(\cdot|x)\|Q_D)=\cR_d(D)
 \label{eq:constant-divergence}
\end{equation}
independent of $x$.  For every input prior $P$, the classical
information-radius identity \cite{CsiszarShields2004} gives
\begin{align}
 I_P(X;Z)
 &=\int\KL(W_D(\cdot|x)\|Q_D)\dd P(x)
   -\KL(PW_D\|Q_D)\\
 &\le\cR_d(D).
\end{align}
Haar input attains equality, proving both capacity and minimax equality for
$0<D<R_{0,d}^2$.  At $D=R_{0,d}^2$ the constant report has zero capacity.
At $D=0$, exact reconstruction of the non-atomic projective source has
infinite capacity, matching $\cR_d(0)=\infty$.
\end{proof}

This specialization is not a novelty claim for invariant, compound-source, or
minimax rate--distortion theory itself; such formulations go back at least to
Sakrison, with compact-group and modern universal-coding developments in
\cite{Sakrison1969,Harremoes2009CompactGroups,MahmoodWagner2022}.  The result
here is the exact evaluation of the stated projective worst-state frontier.

\section{Exact finite-block cost and zero dispersion}

For convenience extend $\cR_d(D)$ by zero for $D\ge R_{0,d}^2$; this
preserves convexity and monotonicity.

Allow an arbitrary joint channel from $n$ pure input states to a shared
classical memory with a standard-Borel output alphabet and $n$ joint reports,
each Borel as a map from the memory into $L^2(\nu)$.  Define
\begin{equation}
 \cC_{d,n}^{\rm univ}(D)=
 \inf_{(W_n,q^n)}\sup_{P_{X^n}} I_P(X^n;Z),
 \label{eq:block-definition}
\end{equation}
where the infimum is subject to
\begin{equation}
 \sup_{x^n}\E\left[
 \frac1n\sum_{t=1}^n\ell_{x_t}(q_{t,Z})\,\middle|\,X^n=x^n
 \right]\le\frac{D}{d(d+1)}.
 \label{eq:block-risk}
\end{equation}

\begin{theorem}[Exact arbitrary-joint-memory finite-block law]
\label{thm:block}
For every integer $n\ge1$ and $0\le D\le R_{0,d}^2$,
\begin{equation}
 \boxed{\cC_{d,n}^{\rm univ}(D)=n\cR_d(D)}.
 \label{eq:block-equality}
\end{equation}
Arbitrary correlations in the representation do not lower the cost; a
memoryless causal product of covariant complex-Bingham channels attains it.
\end{theorem}

\begin{proof}
For the lower bound, choose the iid Haar prior and apply the pointwise
projection coordinatewise.  If $D_t$ is the marginal mean projector
distortion, independence and the mutual-information chain rule give
\begin{equation}
 I(X^n;Z)-\sum_t I(X_t;Z)
 =\sum_t I(X_t;X^{t-1}|Z)\ge0.
\end{equation}
Data processing and the one-letter converse therefore give
\begin{equation}
 I(X^n;Z)\ge\sum_{t=1}^n I(X_t;\widehat\rho_t)
 \ge\sum_{t=1}^n\cR_d(D_t)
 \ge n\cR_d\!\left(\frac1n\sum_tD_t\right),
\end{equation}
where the last step is convexity.  The worst-sequence constraint implies
$n^{-1}\sum_tD_t\le D$ under iid Haar input, so monotonicity yields the
claimed lower bound.  For the upper bound use
$W_D^{\otimes n}$.  Relative to $Q_D^{\otimes n}$, every input sequence has
divergence $n\cR_d(D)$ by \eqref{eq:constant-divergence}; the
information-radius identity upper-bounds capacity by that value, and iid Haar
input attains it.
\end{proof}

We next turn from stochastic representations to fixed-length lossy codes.  An
$(n,M,D,\epsilon)$ code maps an iid Haar source $X^n$ to one of $M$ messages
and then to physical density reports $\widehat\rho^n$, with
\begin{equation}
 \Pp\left\{\frac1n\sum_{t=1}^n
 \|\rho_{X_t}-\widehat\rho_t\|_F^2>D\right\}\le\epsilon.
 \label{eq:excess-distortion-event}
\end{equation}

\begin{theorem}[Zero dispersion and exponential strong converse]
\label{thm:strong-converse}
Let $0<D<R_{0,d}^2$ be a differentiability point of $\cR_d$; this includes
the relative interior of every linear face.  Then the corresponding
$D$-tilted information is constant:
\begin{equation}
 \jmath_X(x,D)=\cR_d(D)
 \quad\text{for Haar-almost every }x.
 \label{eq:constant-jmath}
\end{equation}
Hence the rate dispersion is zero.  Every fixed-length $n$-state code obeys
\begin{equation}
 \boxed{\epsilon\ge
 \left[1-\exp\{\log M-n\cR_d(D)\}\right]_+.}
 \label{eq:strong-converse}
\end{equation}
In particular, if $\log M\le n(\cR_d(D)-\eta)$, then
$\epsilon\ge1-e^{-n\eta}$.
\end{theorem}

\begin{proof}
Compactness of the source and physical reproduction spaces and continuity of
the bounded distortion give an attaining invariant KKT reproduction law
$P^*_{\widehat\rho}$ at every stated $D$.  Set
$\lambda=-\cR_d'(D)>0$ and define
\begin{equation}
 \jmath_X(x,D)=-\lambda D-
 \log\int e^{-\lambda\|\rho_x-\widehat\rho\|_F^2}
 \dd P^*_{\widehat\rho}(\widehat\rho).
 \label{eq:tilted-information-definition}
\end{equation}
On a linear face take the invariant mixture of the supporting active physical
orbit laws.  By transitivity, \eqref{eq:tilted-information-definition} is
independent of $x$; its Haar mean is $\cR_d(D)$ by dual equality.  This
proves \eqref{eq:constant-jmath} and zero variance.  Additivity makes the
$n$-letter tilted information identically $n\cR_d(D)$.  Theorem~7 of Kostina
and Verd\'u \cite{KostinaVerdu2012} states that, for every $\gamma\ge0$,
\begin{equation}
 \epsilon\ge
 \Pp\{\jmath_{X^n}(X^n,nD)\ge\log M+\gamma\}-e^{-\gamma}.
\end{equation}
Taking $\gamma=n\cR_d(D)-\log M$ when positive gives
\eqref{eq:strong-converse}.
\end{proof}

Zero dispersion is a symmetry consequence, not a uniquely quantum effect.
The theorem is a converse; it does not assert a matching constant or
$\tfrac12\log n$ third-order achievability term.

\section{Causal and finite-query consequences}

Exact weighted complex projective $2$-designs replace Haar query integration
in Theorem~\ref{thm:pointwise-projection} without changing any constants.
This makes the probability-field risk finitely witnessable.  Sequential
extensions require an explicit isotropy condition: a policy that always asks
one fixed ray tests only one overlap and cannot force full-state information.

\begin{corollary}[History-dependent conditional-design tensorization]
\label{cor:causal}
Let $X_1,\ldots,X_n$ be fresh independent Haar preparations.  At time $t$, a
causal classical encoder produces $Z_t$ from $(X^t,Z^{t-1})$, and the reporter
may use $Z^t$; any externally randomized menu or seed visible to the reporter
is included in this transcript.  For every allowed history on which the menu may depend,
suppose its conditional scoring measure is a weighted exact projective
$2$-design and the random query draw has no additional access to the fresh
$X_t$.  Let $\widehat\rho_t(Z^t)$ be the measurable physical projection of
the current report and define
\begin{equation}
 D_t=\E\|\rho_{X_t}-\widehat\rho_t(Z^t)\|_F^2.
\end{equation}
Then
\begin{equation}
 I(X^n;Z^n)\ge\sum_{t=1}^n\cR_d(D_t)
 \ge n\cR_d\!\left(\frac1n\sum_tD_t\right).
 \label{eq:causal-bound}
\end{equation}
For the causal factorization
$W(\dd z^n\Vert x^n)=\prod_tW_t(\dd z_t|z^{t-1},x^t)$, define
\begin{equation}
 I(X^n\to Z^n)=\sum_{t=1}^n I(X^t;Z_t|Z^{t-1}).
\end{equation}
For this exogenous iid source, Massey's conservation identity
\cite{Massey1990Directed} gives $I(X^n\to Z^n)=I(X^n;Z^n)$.
\end{corollary}

\begin{proof}
Apply the conditional $2$-design projection and the one-letter converse at
each history.  Independence gives
\begin{align}
 I(X^n;Z^n)
 &=\sum_t I(X_t;Z^n|X^{t-1})\\
 &\ge\sum_t I(X_t;Z^t|X^{t-1})\\
 &=\sum_t I(X_t;Z^t,X^{t-1})\\
 &\ge\sum_t I(X_t;Z^t),
\end{align}
and the one-letter bounds plus convexity complete the proof.  Because the
source law has no feedback dependence on past $Z$, reverse directed
information vanishes.
\end{proof}

At fixed $0<\delta<1$, Theorem~\ref{thm:thermodynamic-rdf} also implies a
classical predictive information cost of
$d\,r_\infty(\delta)/\log2+o(d)$ bits per state.  A direct $d$-level quantum
system has logarithmic Hilbert-space size, but this comparison must not be
misread: a single quantum specimen is not an exact reusable classical
probability oracle, and remote-state preparation or shared entanglement is a
different resource model \cite{BennettRemote2001,BennettEtAlRSP2005}.

\section{Operational boundary and open Grassmann frontier}

The theorems concern probability-field calibration.  They do not derive
Born's rule, simulate state update, certify intrinsic randomness, or exclude
click-only hidden-variable models.  Indeed, reproducing sampled clicks can
require far less information than retaining calibrated probabilities for all
queries.  Channel capacity is an operational asymptotic information resource;
it is not a literal count of physical memory states in one realization.

The natural mathematical extension replaces a ray by a Haar rank-$r$
projector.  Axial overlaps then have beta law $\mathrm{Beta}(r,d-r)$, and the
skewness changes sign under $r\leftrightarrow d-r$.  However, the full
partition is matrix-Jacobi rather than scalar Dirichlet.  An axial ansatz is
not a proof: exact rank-$r$ rate--distortion requires showing that a specified
reconstruction spectrum globally maximizes a matrix-Bingham/HCIZ integral at
fixed Frobenius norm.  Existing complex-projective and Grassmann quantization
results, arbitrary-radius metric-ball asymptotics, and general
rate--distortion bounds do not supply that extremum
\cite{MondalDuttaHeath2006,DaiLiuRider2008,PitavalEtAl2016,RieglerKolianderBolcskei2023}.
We therefore leave it as an explicit open problem.

\section{Conclusion}

We have given a self-contained rate--distortion theory for calibrated Born
prediction on complex projective space.  The all-degree one-spike partition
theorem reduces the unrestricted finite-dimensional problem to an exact
scalar complex-Bingham envelope, including flagged mixtures at coexistence.
It yields a discontinuous directional-information onset for every $d\ge3$,
the exact qubit specialization, and the sharp fixed-$d$ high-fidelity
constant.  In the thermodynamic regime, a uniform projective Laplace principle
gives the closed two-piece frontier, the nontrivial coexistence root
$y-1=2\log y$, and the finite-size correction
$\lambda_{c,d}=\alpha_*d-c_*\log d+O(1)$.  Finally, a pointwise
physical-report projection upgrades Haar-average optimality to exact
worst-state channel capacity, arbitrary-joint-memory tensorization, zero
dispersion, and an exponential strong converse.  The resulting theorem
package connects exact finite-dimensional probability prediction to its
high-dimensional phase law without identifying calibrated prediction with
click generation or particle dynamics.

\section*{Data and code availability}

The accompanying archive \path{repro_v1.zip} is supplied as supplementary
material with this manuscript.  It has the immutable content identifier
SHA-256
\begin{center}
\footnotesize
\texttt{54a0fd17a2b68bdfc62a188043be104ef533b1b20581789a9992ba912902821e}.
\end{center}
From the extracted directory, the exact replay commands are
\begin{verbatim}
python verify_frontier.py
python verify_thermodynamic_limit.py --max-d 100 --output-dir .
\end{verbatim}
They evaluate the finite-dimensional scalar envelope, the coexistence root,
the bounded first-contact problem, and the limiting curve using one process,
without branch-and-bound or heavy search.  The generated JSON files record
the environment, grids, residuals, and runtime.  These outputs are diagnostic;
all headline results are analytic.

\section*{AI assistance and author responsibility}

AI tools assisted with literature search, algebraic stress testing, editing,
and reproducibility checks.  The author is responsible for the definitions,
proofs, claims, citations, and final manuscript.

\appendix

\section{Finite-dimensional spectral reduction}
\label{app:finite}

We give the complete reduction so that every finite-dimensional input used in
the main text is verifiable within this manuscript.

\begin{proposition}[Brazitikos--Pandis normalized interface]
\label{prop:bp-interface}
Let $d\ge2$, let $\gamma_1,\gamma_2,\gamma_3$ be nonnegative integer
multiplicities with $\gamma_1+\gamma_2+\gamma_3=d$, and let $u,v,w\in\R$
satisfy
\begin{equation}
 \gamma_1u+\gamma_2v+\gamma_3w=0,
 \qquad
 \gamma_1u^2+\gamma_2v^2+\gamma_3w^2=1.
 \label{eq:bp-interface-constraints}
\end{equation}
For the grouped vector
\begin{equation}
 x=(\underbrace{u,\ldots,u}_{\gamma_1},
    \underbrace{v,\ldots,v}_{\gamma_2},
    \underbrace{w,\ldots,w}_{\gamma_3}),
\end{equation}
and every positive integer $m$,
\begin{equation}
 p_m(x)\le p_m(a^*_1),
 \qquad
 a^*_1=\left(\sqrt{\frac{d-1}{d}},
 -\frac1{\sqrt{d(d-1)}},\ldots,
 -\frac1{\sqrt{d(d-1)}}\right).
 \label{eq:bp-interface-bound}
\end{equation}
Here zero multiplicities cover the two-value case.  This is precisely
Proposition 5.5 of Brazitikos and Pandis in their unit-norm normalization
\cite{BrazitikosPandis2025}.  Homogeneity gives radius $r$ by multiplying the
right side by $r^m$.
\end{proposition}

\begin{remark}
Proposition~\ref{prop:bp-interface} is the only external extremal inequality
used below.  It asserts the grouped power-sum maximum.  The reduction of an
arbitrary spectrum to at most three values, the odd-sign closure, the Newton
step, the projective Laplace theorem, and the rate--distortion application are
proved here.
\end{remark}

\begin{lemma}[All-degree centered power sums]
\label{lem:appendix-powers}
Let $d\ge3$, $r>0$, and let $a\in\R^d$ satisfy
$\sum_i a_i=0$ and $\sum_i a_i^2=r^2$.  Put
\begin{equation}
 a^*=r\left(\sqrt{\frac{d-1}{d}},
 -\frac1{\sqrt{d(d-1)}},\ldots,
 -\frac1{\sqrt{d(d-1)}}\right).
 \label{eq:appendix-spike}
\end{equation}
Then, for every integer $m\ge2$,
\begin{equation}
 |p_m(a)|\le p_m(a^*),\qquad p_m(a)=\sum_i a_i^m.
 \label{eq:appendix-power-bound}
\end{equation}
\end{lemma}

\begin{proof}
For $m=2$, the assertion is immediate because
$p_2(a)=r^2=p_2(a^*)$ is fixed by the norm constraint.  Hence assume $m\ge3$.
The centered sphere is compact.  Its constraint gradients $\mathbf 1$ and $2a$
are independent because $a$ is nonzero and centered.  At a constrained global
maximum or minimum of $p_m$, Lagrange multipliers give
\begin{equation}
 mx^{m-1}-2\xi x-\zeta=0
 \label{eq:appendix-lagrange}
\end{equation}
for every coordinate value $x$.  The derivative of this sparse polynomial is
$m(m-1)x^{m-2}-2\xi$.  It has at most one real zero when $m$ is odd and at
most two when $m$ is even.  Rolle's theorem therefore leaves at most two,
respectively three, distinct coordinate values at every global extremum.
Proposition~\ref{prop:bp-interface}, applied precisely to this grouped
centered fixed-norm class, places the maximum of $p_m$ at $a^*$.  For odd
$m$, the map $a\mapsto-a$ converts the
minimum to minus the maximum; for even $m$, $p_m(a)\ge0$.  This proves
\eqref{eq:appendix-power-bound}.  The external proposition supplies the
grouped power-sum extremum, not the all-degree Laplace conclusion below.
\end{proof}

If $P\sim\mathrm{Dirichlet}(1,\ldots,1)$, then
\begin{equation}
 \E\langle a,P\rangle^k
 =\frac{k!(d-1)!}{(d+k-1)!}h_k(a),
 \label{eq:dirichlet-hk}
\end{equation}
where $h_k$ is the complete homogeneous symmetric polynomial.  Since
$p_1(a)=0$, Newton's expansion becomes
\begin{equation}
 h_k(a)=
 \sum_{\substack{m_1+2m_2+\cdots+km_k=k\\m_1=0}}
 \prod_{j=2}^k\frac{p_j(a)^{m_j}}{m_j!j^{m_j}}.
 \label{eq:appendix-newton}
\end{equation}
For $d\ge3$, every $p_j(a^*)$ with $j\ge2$ is positive.  Hence
\begin{equation}
 h_k(a)\le |h_k(a)|
 \le\sum_{\substack{m_1+2m_2+\cdots+km_k=k\\m_1=0}}
 \prod_{j=2}^k\frac{|p_j(a)|^{m_j}}{m_j!j^{m_j}}
 \le h_k(a^*).
 \label{eq:appendix-hk-bound}
\end{equation}
Equations~\eqref{eq:dirichlet-hk} and \eqref{eq:appendix-hk-bound}, summed
termwise in the entire exponential series, prove the axial projective
partition inequality in every degree.  When $d=2$, the centered fixed-radius
sphere consists only of the two permutations of
$(r/\sqrt2,-r/\sqrt2)$, so the result is immediate.

It remains to derive and attain the scalar envelope.  Write a density
reproduction as $I/d+A$, where $\Tr A=0$ and $\|A\|_F\le R_{0,d}$.  For
\begin{equation}
 Q_\lambda(A)=\int
 e^{-\lambda\|S_x-A\|_F^2}\dd\mu_d(x),
 \label{eq:appendix-Q}
\end{equation}
conditional Gibbs duality, applied to every posterior given $A$, gives
\begin{equation}
 I(X;A)\ge-\lambda D-\E\log Q_\lambda(A)
 \ge-\lambda D-\log\sup_AQ_\lambda(A).
 \label{eq:appendix-gibbs}
\end{equation}
At fixed $\|A\|_F$, diagonalization and the all-degree partition inequality
replace its spectrum by the valid axial density comparator
$A=b(\rho_u-I/d)$, $0\le b\le1$.  For this comparator,
\begin{equation}
 \log Q_\lambda(A)
 =-\lambda R_{0,d}^2+K_d(2\lambda b)
  -\lambda R_{0,d}^2b^2
 =-\lambda R_{0,d}^2+G_{d,\lambda}(b).
 \label{eq:appendix-axial-Q}
\end{equation}
Maximizing first over $b$ and then over $\lambda$ proves the converse in
Theorem~\ref{thm:finite-envelope}.

For equality, let $b>0$ be active, put $\kappa=2\lambda b$, and take
\begin{equation}
 \frac{\dd P_{X|U=u}}{\dd\mu_d}(x)
 =\frac{e^{\kappa|\langle x,u\rangle|^2}}{M_d(\kappa)},
 \qquad U\sim\mu_d.
 \label{eq:appendix-bingham}
\end{equation}
Invariance makes the source marginal Haar.  If
$m_d(\kappa)=\partial_\kappa\log M_d(\kappa)$, stationarity of the active
radius is
\begin{equation}
 b=\frac{d m_d(\kappa)-1}{d-1}.
 \label{eq:appendix-self-consistency}
\end{equation}
Thus the conditional mean is exactly $(1-b)I/d+b\rho_u$, so the channel is
its own posterior-mean reproduction.  It attains
\begin{equation}
 D_b=R_{0,d}^2(1-b^2),\qquad
 I_b=\kappa m_d(\kappa)-\log M_d(\kappa)
 =\lambda(R_{0,d}^2-D_b)-\max_{0\le v\le1}G_{d,\lambda}(v).
 \label{eq:appendix-branch}
\end{equation}
The optimized log partition
$c_d(\lambda)=-\lambda R_{0,d}^2+\max_bG_{d,\lambda}(b)$ is convex, and
Danskin's theorem gives
\begin{equation}
 -\partial c_d(\lambda)
 =\operatorname{conv}\{D_b:b\text{ is active at }\lambda\}.
 \label{eq:appendix-danskin}
\end{equation}
Independently flagged mixtures of the corresponding invariant channels attain
every distortion in this interval.  The independent channel gives
$D=R_{0,d}^2$, $I=0$; exact reconstruction is obtained as the limiting
endpoint and has infinite information.  Thus all
$0\le D\le R_{0,d}^2$ are covered without assuming uniqueness of the
finite-dimensional radial maximizer.

\printbibliography

\end{document}
